\chapter{}

Quel poco tempo che Anna stava al podere, quando non aveva più da lavorare in casa, si faceva empire le brocche da Ghìsola; e poi, con un annaffiatoio, bagnava le piante dei limoni. La sera Giacco toglieva, con una zappa, l'erba nata attorno alla casa; buttandola ai conigli o alle galline.

Anna scendeva fin giù agli orti, e qualcuna delle donne le pigliava l'insalata e i cavoli.

Ella avrebbe voluto tenere i fiori, anche perchè vicino a Poggio a' Meli c'era un giardino; che andava sempre a vedere, per ambire di averne un altro eguale. Ma dovevano bastarle i gerani e i garofani; quando glieli regalavano da trapiantare. Non osava, però, tenerne molti, perchè certo Domenico le avrebbe domandato se andava in campagna per curarsi oppure per starci in villeggiatura. Del resto ella stessa si contentava d'averne più di quando era ragazza.

Anche per comprare quei pochi ninnoli che teneva nel suo salotto di città, era bisognato che glieli avessero venduti quasi per forza. Infatti un ebreo robivendolo, tutte le volte che non aveva da pagare il conto alla trattoria, le portava a far vedere ogni specie di oggetti vecchi e glieli lasciava sul banco; benché lei non volesse in nessun modo. E quando, passata una settimana, egli tornava, Domenico ed Anna, dopo mezz'ora per mettersi d'accordo, e avergli detto che sarebbe stata l'ultima volta, si aggiustavano alla meglio. Il robivendolo giurava che da qui in avanti avrebbe pagato sempre con i soldi alla mano; e allora bevevano insieme un bicchiere di vino, perchè erano doventati anche rochi a forza di vociare e di trattarsi male.

Ma Anna ne era contenta; e così i quadri, dipinti sul vetro, delle Cinque parti del mondo, i portafiori d'alabastro ingiallito, le anfore di vera porcellana entravano in casa sua.

Il salotto, ormai, non ne conteneva più. C'era poi addirittura una parete ricoperta con le fotografie di quasi tutti i conoscenti; e, sopra un mobile verniciato a noce, due ciociare di gesso che sorridevano. Nel tavolino di mezzo, un servito di cristallo celeste, ma incompleto; che aveva attorno cinque lucernine di ottone sempre infioccate su nel manico perchè le mandavano, con un fiasco d'olio, a tenerle accese quando facevano i Sepolcri.

Ella dava, almeno una volta al mese, il cinabro agli impiantiti; e, allora, bisognava che si pulissero bene le scarpe prima d'entrare.

Quando, in campagna, le portavano qualche fiore, non voleva tenerlo in casa; e l'offriva alla Madonna del Convento di Poggio al Vento. Se fosse stato già tardi e avevano chiusa la chiesa, lo metteva in fresco, ma sopra il tavolo della stanza d'ingresso; e la mattina dopo era la prima faccenda.

Per pararsi il sole, che le faceva subito dolere la testa, aveva un ombrellino rosso con il manico d'avorio; un ombrellino di parecchi anni. Ella, quando vedeva le assalariate, se ne vergognava; e, chiudendolo, stava piuttosto sotto una pianta. Mentre invece, andando alla messa, lo portava volentieri; e magari se lo faceva reggere da Ghìsola.

In chiesa si metteva su una panca, un poco distante dalle contadine; che, del resto, per rispetto, a farle posto ci pensavano anche da sé.

S'era fatto un vestito nero con una guarnizione di seta gialla al collo; e con una trina che, attaccata alle spalle e alla cintura, stava fino a mezze maniche. Su la guarnizione teneva una catena d'oro. Invece, per la trattoria, aveva un vestito rosso a palline bianche e celesti.

Ella diceva a Ghìsola che imparasse a scrivere, almeno un poco; ma siccome non poteva fidarsi che Pietro le insegnasse, perchè si metteva subito a farle dispetti, lei stessa ci si dedicava qualche ora del giorno, quando stava meglio. E Ghìsola s'era fatto l'inchiostro con le more delle siepi. Ma non andò mai avanti oltre le prime aste.

Per dire la verità, invece, Ghìsola avrebbe imparato volentieri, e, a sapere che Pietro andava a scuola, le faceva un grande effetto. Ella avrebbe voluto almeno leggere, perchè molte delle sue amiche dei poderi accanto avevano perfino il libro da messa, quello regalato dai cappuccini per la prima comunione; e poi perchè in Piazza del Campo, le domeniche mattine, le veniva voglia di comprare le canzonette stampate che vendevano a un soldo con il racconto di qualche fatto miracoloso, dove c'era sempre una Madonna con una gran corona dietro la testa. Le canzonette erano belle perchè anch'esse, prima delle rime, ci avevano sempre qualche figurina. Ella si fermava, con gli altri contadini, a sentirle con la chitarra da Cicciosodo, quel cantastorie capace di smuovere il cappello a tuba contraendo la pelle della fronte, ritto sopra uno sgabello. C'erano anche le scimmie che sceglievano i numeri con la Ruota della Fortuna; e c'era chi vendeva certe chicche di tutti i colori, involtate in cartocci ritagliati a frange con le forbici.

Quando tornava a Poggio a' Meli, aveva già imparato l'aria della canzonetta che le era piaciuta di più; ma non si ricordava di tutte le parole. Qualche volta, avendola comprata lo stesso e tenendola piegata in tasca perchè Masa non gliela vedesse, se le faceva leggere quando nel campo trovava un'amica. C'erano da vero cose belle, che la commovevano o la facevano ridere.
